\documentclass[a4paper,12pt]{article}

\usepackage{graphicx}
\usepackage{hyperref}
\usepackage[italian]{babel}
\usepackage{parskip}
\usepackage{multicol}
\usepackage{verbatim}
\usepackage[a4paper,left=2cm,right=2cm,top=2.5cm,bottom=2.5cm]{geometry}

\setlength{\parskip}{1em}
\graphicspath{ {./images/} }
\hypersetup {
    colorlinks=true,
    linkcolor=black,
    pdfborderstyle={/S/U/W 1},
}

\begin{document}
\begin{titlepage}
    \begin{center}
        \begin{minipage}{0.3\textwidth}
            \includegraphics[width=\linewidth]{Logo_Universita_Padova300x300.png}
        \end{minipage}
        
    \end{center}

    \begin{center}
        \noindent\rule{0.7\textwidth}{0.4pt}\\[0.5cm]
        \large\textbf{Università degli Studi di Padova}\\
        \large Corso di Laurea in Informatica\\
        \small A. A. 2025/26
    \end{center}

    \vspace{1,5cm}    

    \begin{center}
        \begin{minipage}{0.5\textwidth}
            \includegraphics[width=\linewidth]{logo.png}
        \end{minipage}
    \end{center}

    \vspace{1,5cm}    
    \begin{center}
        {\Large \textbf{Relazione}\par}
        \vspace{0.6cm}
        {\Large Progetto di Tecnologie Web\par}
        \vspace{0.3cm}
    \end{center}
    \vfill

    \noindent
    \begin{tabular}{r|l}
        \textbf{Studenti} & Sofia De Blasi - matricola - email \\
         & Nenad Radulovic - matricola - email \\
         & Giacomo Speggiorin - matricola - email \\
         & Nicola Simionato - matricola - email \\
         & \\
        \textbf{Indirizzo del sito} & indirizzo\\
         & \\
        \textbf{Utenti} & 
        \textbf{Username:} user, \textbf{Password:} user \\ 
        & \textbf{Username:} admin, \textbf{Password:} admin\\


    \end{tabular}

    \vfill

\end{titlepage}     
\newpage
\tableofcontents
\newpage

\section{Introduzione}
    \subsection{Com'è nata l'idea}
L'idea di questo progetto nasce da una domanda che ci siamo posti più volte nella nostra quotidianità: \textit{"Cosa potremmo fare questo weekend?"}. 
Nelle città universitarie vengono spesso organizzati eventi e attività dedicate agli studenti, è però seccante doverne andare ogni volta alla ricerca. 
Bisogna sapere innanzitutto cosa si sta cercando (il che non è una cosa scontata), e serve poi andarsi a sfogliare di volta in volta le pagine social degli organizzatori, che sono sempre diversi.\par
Abbiamo quindi pensato a quanto sarebbe comodo e utile se fossero gli organizzatori stessi ad andare incontro agli studenti, e a quanto si semplificherebbe la ricerca se esistesse un'unica bacheca dove tutti gli organizzatori possano appendere i loro volantini.\par
Nel mentre buttavamo giù le idee il discorso si è ampliato, e ci siamo resi conto che la fuori ci sono un sacco di servizi di cui gli studenti hanno bisogno ma non hanno un punto di riferimento sul dove andarli a cercare.

    \subsection{L'obbiettivo di StudentSpace}
StudentSpace propone una \textit{bacheca online} che raccolga in un unico luogo eventi, attività e servizi di interesse per gli studenti, andando oltre il contesto strettamente accademico.\par
Attraverso la piattaforma sarà possibile trovare eventi sportivi, attività serali, concerti, e altre occasioni di socialità.
Potranno inoltre dare la loro disponibilità ad altri studenti per ripetizioni, oppure consultare annunci di affitto relativi alla propria città.\par
L'obbiettivo principale è garantire a studenti e studentesse un accesso rapido e completo a ciò che potrebbe interessarli.
Allo stesso tempo, si vuole agevolare gli organizzatori, che si rivolgono a questo target, fornendo loro un canale unico e mirato per raggiungere un pubblico ampio e specifico.

\section{Analisi del lavoro svolto}
    \subsection{Processo decisionale}
Durante tutto lo sviluppo del progetto ci si è trovati a dover prendere diverse decisioni in merito a design, funzionalità, contenuto ecc. 
Nella maggior parte dei casi questi punti suono stati discussi in presenza, tra tutti i membri del gruppo. 
Ovviamente, questo metodo comporta delle chiare limitazioni, si è quindi deciso di adottare
anche dei metodi di comunicazione (sia sincroni che asincroni) alternativi.\\\\
Per quanto riguarda le comunicazioni asincrone, è stata usata ampiamente la piattaforma di messaggistica \textbf{\underline{\nameref{strumenti:whatsapp}}}. 
Si è deciso inoltre di adottare \textbf{\underline{\nameref{strumenti:discord}}} come strumento di comunicazione asincrona, in supporto alle riunioni in presenza.

    \subsection{Modalità di lavoro}
Per adattarsi meglio alle esigenze del gruppo si è deciso che i membri del gruppo avrebbero svolto ciascuno le proprie attività in maniera prevalentamente asincrona.
Questa decisione è stata presa per riuscire a dare più libertà ai singoli membri nell'organizzare il proprio lavoro. 
La scelta è stata inoltre supportata dalla volontà del gruppo di dare maggiore fluidità allo sviluppo del progetto. 
Il lavoro è stato infatti suddiviso in modo che ogni membro andasse a lavorare su punti diversi rispetto agli altri, cosicché nessuno dovesse preoccuparsi che il proprio lavoro andasse in conflitto con quello di qualcun altro. 
A sostegno di ciò si è deciso di implementare un sistema di versionamento del codice (con annesso hosting online) tramite \textbf{\underline{\nameref{strumenti:git}}} e \textbf{\underline{\nameref{strumenti:github}}}.

    \subsection{Divisione del lavoro}
Visto lo scopo di natura accademica del progetto si è ritenuto opportuno che ciascun membro del gruppo fosse coinvolto (quantomeno in parte) in ogni ambito del progetto.
Questo ha permesso a ogni membro di avere un idea generale sul funzionamento del programma, e di mettere mano direttamente a ogni fase del suo sviluppo.

\section{Analisi dei requisiti}
    \subsection{Tipologie di Utenti}
        \begin{itemize}
            \item \textbf{Utente non registrato:} utente che esplora le categorie, legge annunci e capisce  come funziona la bacheca
                \begin{itemize}
                    \item  \textbf{Azioni e permessi:}  vidualizzazione completa delle varie pagine, navigazione tra categorie e annunci, ricerca e filtraggio, apertura dettaglio annuncio, avvio del flusso di registrazione/login
                    \item \textbf{Limitazioni:} accesso negato alle funzioni di pubblicazione e gestione $\rightarrow$  inviti chiari a registrarsi quando tenta azioni riservate.
                \end{itemize}
            \item \textbf{Utente registrato (semplice):} utente che oltre alle funzionalotà che erdità dall'untente non registrato può pubblicare annunci e gestirli.
                \begin{itemize}
                    \item \textbf{Azioni e permessi:}
                    creazione, modifica ed eliminazione dei propri annunci; gestione profilo; compilazione form con parte generica e specifica, revisione
                    \item \textbf{Limitazioni:} non può modificare annunci altrui; vincoli su campi obbligatori e formati (validazione client/server).
                \end{itemize}
            \item \textbf{Amministratore: utente che modera e garantisce il giusto funzionamento} % da decidere
                \begin{itemize}
                    \item \textbf{Azioni e permessi:}gestione completa degli annunci (anche altrui), categorie, utenti; strumenti di moderazione (blocco, archiviazione)
                \end{itemize}
        \end{itemize}

    \subsection{Analisi dei contenuti} % Tipi di annunci, Campi generici, Campi specifici per categoria
    Il sito web si configura come una bacheca universitaria online, in cui gli utenti possono
    pubblicare e consultare annunci suddivisi in quattro categorie principali: \textbf{Affitti},
    \textbf{Esperimenti}, \textbf{Eventi} e \textbf{Ripetizioni}. L’analisi dei contenuti si concentra sui tipi di annunci, sui campi comuni a tutte le categorie e sui campi specifici per ciascuna categoria.
    
        \subsubsection{Tipi di annunci}
        \begin{itemize}
          \item \textbf{Affitti}: annunci relativi a stanze, appartamenti o posti letto disponibili per studenti.
          \item \textbf{Esperimenti}: annunci riguardanti attività di laboratorio, progetti di ricerca o collaborazioni scientifiche.
          \item \textbf{Eventi}: annunci di conferenze, seminari, feste universitarie, attività culturali o sportive.
          \item \textbf{Ripetizioni}: annunci di studenti che offrono supporto didattico su materie specifiche.
        \end{itemize}

        \subsubsection{Campi generici}
        Ogni annuncio, indipendentemente dalla categoria, deve contenere i seguenti campi:
        \begin{itemize}
          \item \textbf{Categoria}: selezionata dall’utente tra le quattro disponibili.
          \item \textbf{Titolo}: breve descrizione sintetica dell’annuncio.
          \item \textbf{Data di pubblicazione}: generata automaticamente dal sistema.
          \item \textbf{Città}: riferimento all’utente registrato che ha pubblicato l’annuncio.
          \item \textbf{Descrizione}: testo libero con i dettagli dell’annuncio. % decidere
          \item \textbf{Autore}: riferimento all’utente registrato che ha pubblicato l’annuncio. % decidere
        \end{itemize}

        \subsubsection{Campi specifici per categoria}
        Oltre ai campi generici, ciascuna categoria prevede attributi specifici:
        \begin{itemize}
            \item \textbf{Affitti}:
            \begin{multicols}{2}
                \begin{itemize} % decidere se tenere tutti questi attrubuti
                    \item Prezzo mensile
                    \item Indirizzo
                    \item Contatti
                    \item Numero inquilini
                    \item Tipologia
                    \item Piano
                    \item Piani edificio
                    \item Ascensore
                    \item Superficie
                    \item Locali
                    \item Camere da letto
                    \item Bagni
                    \item Arredato
                    \item Balcone
                    \item Terrazzo
               \end{itemize}
           \end{multicols}
            \item \textbf{Esperimenti}:
            \begin{itemize}
              \item Laboratorio di riferimento
              \item Durata prevista
              \item Compenso
            \end{itemize}
          \item \textbf{Eventi}:
            \begin{itemize}
              \item Data evento
              \item Luogo
              \item Tipologia evento (seminario, festa, conferenza, ecc.)
            \end{itemize}
          \item \textbf{Ripetizioni}:
            \begin{itemize}
              \item Materia
              \item Livello (scuola superiore, università, ecc.)
              \item Prezzo orario
            \end{itemize}
        \end{itemize}

\subsection{Requisiti funzionalità} % registrazione/login, pubblicazione annunci, modifica/eliminazione, filtraggio annunci, validazione input (client/server), accessibilità e responsività
Il sistema deve garantire una serie di funzionalità fondamentali per soddisfare i requisiti del progetto
e le esigenze degli utenti.

\subsubsection{Registrazione e login}
\begin{itemize}
  \item Form di registrazione con campi obbligatori (nome, email, password).
  \item Validazione lato client (formato email, lunghezza password) e lato server.
  \item Login con credenziali; gestione sessione sicura.
  \item Utente amministratore con credenziali \texttt{admin/admin}, utente semplice con \texttt{user/user}.
\end{itemize}

\subsubsection{Pubblicazione annunci}
\begin{itemize}
  \item Form di inserimento con parte generica (titolo, descrizione, categoria).
  \item Campi specifici mostrati dinamicamente in base alla categoria scelta.
  \item Salvataggio dati nel database con associazione all’utente autore.
  \item Possibilità di visualizzare un’anteprima prima della conferma. % da decidere
\end{itemize}

\subsubsection{Modifica ed eliminazione annunci}  % da decidere
\begin{itemize}
  \item Utente registrato può modificare o cancellare solo i propri annunci.
  \item Amministratore può intervenire su qualsiasi annuncio.
\end{itemize}

\subsubsection{Filtraggio e ricerca annunci}
\begin{itemize}
  \item Filtri generici (categoria, data di pubblicazione).
  \item Filtri specifici 
  \item Ricerca testuale su titolo e descrizione.
\end{itemize}

\subsubsection{Validazione input}
\begin{itemize}
  \item \textbf{Lato client}: controlli immediati con JavaScript (campi obbligatori, formati corretti).
  \item \textbf{Lato server}: verifica dati ricevuti prima dell’inserimento nel database.
  \item Messaggi di errore chiari e accessibili.
\end{itemize}

\subsubsection{Accessibilità e responsività}
\begin{itemize}
  \item Layout realizzato con CSS3, Flexbox e Grid.\
  \item Separazione netta tra contenuto (HTML5), presentazione (CSS) e comportamento (JS/PHP).
  \item Compatibilità cross-browser e adattamento a diverse dimensioni di schermo.
  \item Uso di etichette semantiche, contrasto cromatico adeguato e font leggibili.
\end{itemize}


\section{Progettazione}
    \subsection{Archietettura del sistema} % struttura delle cartelle, separazione contenuto/presentazione/comportamento, panoramica componenti

\subsection{Progettazione Database} 
\subsubsection{Analisi dei requisiti del database}

La base di dati deve supportare la gestione degli utenti e degli annunci pubblicati
sulla bacheca universitaria. In particolare, deve garantire la memorizzazione dei
campi generici comuni a tutte le categorie di annunci e dei campi specifici per ciascuna
categoria. Inoltre, deve rispettare i vincoli di integrità e consentire operazioni di
inserimento, modifica e cancellazione.

\paragraph{Entità principali}
\begin{itemize}
  \item \textbf{Utente}: rappresenta gli utenti registrati al sistema.
    \begin{itemize}
      \item Attributi: \texttt{IdUtente}, \texttt{Nome}, \texttt{Cognome}, \texttt{Città} \texttt{Email}, \texttt{Password}, \texttt{Ruolo}.
    \end{itemize}
  \item \textbf{Annuncio}: rappresenta un contenuto pubblicato sulla bacheca.
    \begin{itemize}
      \item Attributi: \texttt{IdAnnuncio}, \texttt{Titolo}, \texttt{Descrizione}, \texttt{DataPubblicazione}, \texttt{Città}, \texttt{Categoria}, \texttt{IdUtente}.
    \end{itemize}
  \item \textbf{Categorie di annunci}: specializzazioni dell’entità \texttt{Annuncio}.
    \begin{itemize}
      \item \textbf{Affitto}: \texttt{PrezzoMensile}, \texttt{Indirizzo}, \texttt{Contatti}, \texttt{NumeroInquilini}, \texttt{Tipologia}, \texttt{Piano}, \texttt{PianiEdificio}, \texttt{Ascensore}, \texttt{Superficie}, \texttt{Locali}, \texttt{CamereLetto}, \texttt{Bagni}, \texttt{Arredato}, \texttt{Balcone}, \texttt{Terrazzo}.
      \item \textbf{Esperimento}: \texttt{Laboratorio}, \texttt{DurataPrevista}, \texttt{Compenso}.
      \item \textbf{Evento}: \texttt{DataEvento}, \texttt{Luogo}, \texttt{TipologiaEvento}.
      \item \textbf{Ripetizione}: \texttt{Materia}, \texttt{Livello}, \texttt{PrezzoOrario}.
    \end{itemize}
\end{itemize}

\paragraph{Relazioni}
\begin{itemize}
  \item Un \textbf{Utente} può pubblicare da zero a molti \textbf{Annunci} (relazione 1:N).
  \item Ogni \textbf{Annuncio} appartiene ad una sola \textbf{Categoria} (relazione 1:1 con la specializzazione).
  \item Ogni categoria eredita i campi generici da \texttt{Annuncio} e aggiunge i propri campi specifici.
\end{itemize}

\paragraph{Generalizzazioni}
L’entità \texttt{Annuncio} funge da \textbf{generalizzazione} per le quattro categorie (\texttt{Affitto}, \texttt{Esperimento}, \texttt{Evento}, \texttt{Ripetizione}).
Questa scelta progettuale consente di mantenere un modello chiaro e normalizzato, evitando ridondanze e facilitando la gestione dei dati.

\subsubsection{Schema ER}
\begin{figure}[h] 
    \centering
    \includegraphics[width=1\textwidth]{ER.png}
    \caption{Schema ER}
    \label{fig:ER}
\end{figure}


\subsubsection{Generalizzazioni e sceltre progettuali}
\subsubsection{Schema logico (relazionale)}
\subsubsection{Normalizzazione}
\subsubsection{Vincoli di integrità}


\subsection{Struttura del sito} %mappa delle pagine, funzioni, URL
\subsection{Mockup}
\subsection{Responsive Web Design e Mobile First} 

\section{Sviluppo} % decidere struttura
\begin{comment}
\subsection{HTML} % sementica
\subsection{CSS}
\subsection{JavaScript} % validazione client + interazione
\subsection{PHP}
\subsection{Database}
\subsection{Errori}
\subsection{Validazione} 
\subsection{Gestione delle immagini} 
\subsection{Palette colori e font} 
\subsection{Configurazione .htaccess} % password hashing, SQL injection, sanitizzazione input, protezione aree riservate, gestione sessioni
\end{comment}
% o struttura più sintetica

\subsection{Front-end (HTML, CSS, JS)}
\subsection{Back-end (PHP, Database)}
\subsection{Validazione e Sicurezza}
\subsection{Design (immagini, palette colori)} 

\section{Test}
\subsection{Test funzionali}
Descrizione dei test su login, registrazione, pubblicazione e filtraggio annunci.

\subsection{Test di compatibilità}
Test su diversi browser (Chrome, Firefox, Edge, Safari) e dispositivi (desktop, tablet, smartphone).

\subsection{Test di accessibilità}
Color contrast, navigazione da tastiera, screen reader.

\section{Accessibilità}
\subsection{Introduzione} 
\subsection{WAI-ARIA e screen-reader} 
\subsection{Progressive Enhancement} 
\subsection{Supporto ai dispositivi “legacy”} 
\subsection{Ulteriori funzionalità di accessibilità}

\section{SEO e prestazioni} %meta tag, URL parlanti, sitemap, compressione immagini, lazy loading, minificazione
\subsection{Analisi delle possibili ricerche su Google} % “affitto stanza Padova”, “eventi Padova weekend”, “ripetizioni informatica Padova”…)

\section{Strumenti utilizzati}
\begin{comment}
   \subsection{Canva} \label{strumenti:canva}
\underline{\url{https://www.canva.com/it_it/}} \\\\
Piattaforma di progettazione grafica per creare contenuti visuali per presentazioni, siti web e prodotti simili. 

\subsection{Discord} \label{strumenti:discord}
\underline{\url{https://discord.com/}} \\\\
Piattaforma di messaggistica istantanea, di distribuzione digitale e di voice over IP, utilizzata per la comunicazione all'interno server Discord.
Gli utenti comunicano con chiamate vocali, videochiamate, messaggi di testo, media e file in chat private o come membri di un server Discord. Quest'ultimi sono una raccolta di canali di tipo vocale e/o testuale. 

\subsection{Git} \label{strumenti:git}
\underline{\url{https://git-scm.com/}} \\\\
Applicazione software per il controllo di versione distribuito utilizzabile da interfaccia a riga di comando. Viene utilizzato il versionamento delle repository dedicate al codice e alla documentazione del progetto.

\subsection{Github} \label{strumenti:github}
\underline{\url{https://git-scm.com/}} \\\\
È un servizio di hosting per progetti software che implementa lo strumento di controllo versione distribuito Git.  

\subsection{LaTeX} \label{strumenti:latex}
\underline{\url{https://www.latex-project.org/get/}} \\\\
Linguaggio di marcatura del testo, usato ampiamente in ambito accademico per la preparazione di testi.

\subsection{Silktide} \label{strumenti:silktide}
\underline{\url{https://silktide.com/}} \\\\
È un'estensione che facilita i test sull'accessibilità.

\subsection{W3C Validator} \label{strumenti:w3cvalidator}
\underline{\url{http://validator.w3.org/}} \\\\
Validatore markup HTML/XHTML.

\subsection{WhatsApp} \label{strumenti:whatsapp}
\underline{\url{https://www.whatsapp.com/}} \\\\

\subsection{WebAIM} \label{strumenti:webaim}
\underline{\url{https://webaim.org/resources/contrastchecker/}} \\\\
\end{comment}

\begin{itemize}
  \item \textbf{Canva}  
  \underline{\url{https://www.canva.com/it_it/}} \\
  Piattaforma di progettazione grafica per creare contenuti visuali per presentazioni, siti web e prodotti simili.

  \item \textbf{Discord}  
  \underline{\url{https://discord.com/}} \\
  Piattaforma di messaggistica istantanea e distribuzione digitale, utilizzata per la comunicazione del gruppo. 
  Supporta chiamate vocali, videochiamate, messaggi di testo e condivisione file.

  \item \textbf{Git}  
  \underline{\url{https://git-scm.com/}} \\
  Applicazione software per il controllo di versione distribuito, usata per il versionamento del codice e della documentazione.

  \item \textbf{Github}  
  \underline{\url{https://github.com/}} \\
  Servizio di hosting per progetti software che implementa Git, utilizzato per la collaborazione e la condivisione del codice.

  \item \textbf{LaTeX}  
  \underline{\url{https://www.latex-project.org/get/}} \\
  Linguaggio di marcatura del testo, usato per la preparazione della relazione accademica.

  \item \textbf{Silktide}  
  \underline{\url{https://silktide.com/}} \\
  Estensione che facilita i test sull’accessibilità.

  \item \textbf{W3C Validator}  
  \underline{\url{http://validator.w3.org/}} \\
  Strumento per validare il markup HTML/XHTML.

  \item \textbf{WhatsApp}  
  \underline{\url{https://www.whatsapp.com/}} \\
  Piattaforma di messaggistica istantanea usata per la comunicazione rapida tra i membri del gruppo.

  \item \textbf{WebAIM}  
  \underline{\url{https://webaim.org/resources/contrastchecker/}} \\
  Strumento online per verificare il contrasto dei colori e migliorare l’accessibilità.
\end{itemize}


\section{Guida all’installazione del progetto}


\end{document}

